\appendix
\section{Proof scope and supporting artifacts}\label{sec:scope}
The proof of Theorem~\ref{thm:main} covers every $g\geq2$ and $r\geq2$
with the quantifiers stated there. Rank two is imported with explicit
credit; for $r\geq3$ Sections~\ref{sec:arithmetic}--\ref{sec:closure}
give the construction and proof. Classical surface-character geometry and
Goldman's Hamiltonian theorem are cited inputs. The irreducibility input
is the determinant-one intermediate assertion of the pinned manuscript,
identified in Section~\ref{subsec:geometry}. The proof neither classifies
all arithmetic orbits nor asserts analytic density of the selected orbit.

\subsection{Limited algebraic formal support}
A separate Lean artifact, \texttt{t027o-formal-v1}, supports a precisely
bounded part of Lemma~\ref{lem:minors}. Its theorem
\texttt{triangle\_coboundary}, over an arbitrary field and an index type
with a chosen reference index, proves that a zero-diagonal antisymmetric
array with vanishing oriented triangle sums is of the form
$x_{ij}=h_i-h_j$. Its theorem \texttt{principal\_minor\_rigidity} assumes
nonzero off-diagonal entries, zero diagonal variation, the explicit pair
stationarity polynomials, the six-term cycle derivative polynomials, and
unequal forward and reverse triangle products. It concludes
$E_{ij}=(h_i-h_j)M_{ij}$.

Those formal hypotheses are scalar polynomial identities. The pair identity
becomes the negative two-by-two minor derivative after the diagonal
identities; the cycle identity becomes the three-by-three minor derivative
only after the diagonal and pair identities. The artifact
does \emph{not} prove that they are the derivatives of the corresponding
determinant functions. It does not formalize generic Cartan choices,
Vandermonde interpolation, the surface curves, the residual degeneration,
arithmetic integrality, torus density, Goldman geometry, the analytic
ideal argument, or orbit density. It is not a formal proof of
Theorem~\ref{thm:main}. Its controls include a generic rational example
and a counterexample to deleting the triangle-genericity hypothesis.
This counterexample shows that genericity cannot simply be deleted; it does
not prove necessity of every triple condition in larger index sets. An
independent review adds a formal derivative bridge and further controls in
its own file, \texttt{T027OReview.lean}; those additions are distinct from
the author package and still do not cover the geometric orbit argument.
The exact coverage and trusted build boundary are documented in its
\texttt{t027o-formal-v1/README.md};
the manuscript packet supplies an external-path crosswalk for local use.
Formal review and release status are kept in the project records, not
inferred from the existence of this support file.

\subsection{Attribution and reproducibility limits}
The arithmetic unit construction, principal-minor cycle identities,
Vandermonde arguments, torus saturation, Baire theorem, coherent-ideal
argument, and inverse function theorem are supporting tools. They are not
separate novelty claims. The higher-rank contribution presented here is
the connected construction consisting of a single arithmetic seed,
the two-band tangent certificate at its complex torus family, and the
arithmetic-first, complex-second parameter choice. The literature
comparison in Section~\ref{sec:context} is bounded, with no claim of
complete priority clearance.

The source revision used is
\texttt{fd4aeeb2ee4fc729c18d98444fed42fd0529eeeb} of
\href{https://github.com/openai/math}{OpenAI's math repository}.
The author packet records exact source identities and reading limits.
Private copies of third-party bodies are excluded from the public source
bundle. The accompanying build script recompiles this article from a
copied source directory and records source and PDF hashes. These hashes
and finite diagnostic computations establish identity or finite checks;
they do not replace the mathematical proof.

\paragraph{Acknowledgment of assistance.}
The argument was developed with AI-assisted calculations, source
comparison, and independent project review. Author-side checks included
the unit-independence argument, the based-word conventions, the second-band
derivative, the residual symmetric/skew witness, and the spectral-cover
closure argument. Responsibility for the exposition and its stated
limitations remains with the named author. No institutional affiliation
is asserted.
